

\addcontentsline{toc}{section}{Заключение}
\section*{Заключение}
В результате выполнения работы была создана база данных для предметной области <<Индивидуальный учет поголовья коров>>.

На первом этапе была выбрана и проанализирована предметная область <<Индивидуальный учет поголовья коров>>. Были сформулированы 4 цели функционирования базы данных, после чего было выделено 11 сущностей: заболевания, лактация, порода, ветеринар, дояр, зоотехник, корма, рацион, порция, журнал, корова. Была построена ER-диаграмма, затем --- схема связи объектов, при составлении которой количество сущностей увеличилось до 13, так как были выделены дополнительные словари: период, вид корма. Также на схеме связи объектов указаны связи между объектами: 8 связей типа <<один к многим>>, 5 связей --- <<многие ко многим>>.

При выполнении второго этапа итоговое число сущностей составило 12. Такие сущности, как заболевания, ветеринар, дояр, журнал были удалены из схемы базы данных, сущность период была разделена на две: суточность и сезонность. После составления таблиц с описанием таблиц базы данных, где были аргументированы 13 ограничений на размеры атрибутов, и лингвистического описания было составлено графическое описаине схемы базы данных. На графическом описании схемы базы данных 12 таблиц, над каждой таблицей было подписано количество записей, правила генерации и уровни заполнения. Всего получилось 4 уровня заполнения и 4 правила заполнения. Затем был написан скрипт на языке SQL для генерации таблиц (СУБД PostgreSQL, версия: psql (PostgreSQL) 14.22). Для заполнения таблиц первого уровня был также написан скрипт на языке SQL, а заполнение таблиц 2-4 уровней было реализовано с помощью языка Python (версия: Python 3.10.12) и библиотеки faker. После заполнения таблиц общее количество записаей составило 302 526, где количетсво записей в таблице Breed --- 10, в таблице Day\_Times --- 5, в Season --- 4, в Lactation\_Phase --- 5, в Feed\_Type --- 17, в Livestock\_Specialist --- 8, в Cow --- 6574, в Feeds --- 40, в Diet --- 38, в Lactation --- 130910, в Feeding\_Plan --- 464, в Nutrition --- 164451.

На третьем этапе были сформулированы 9 запросов к базе данных на языке SQL, объединившие в себе такие логические группы инструментов, как операторы манипулирования данными (SELECT, UPDATE, SET), средства соединения таблиц (JOIN, LEFT JOIN), предикаты фильтрации и группировки (WHERE, HAVING, AND, GROUP BY), агрегатные функции (COUNT, MAX, MIN), а также модификаторы выборки и псевдонимы (DISTINCT, AS). Затем были построены гистограммы: две 2d-гистограммы для запросов 3 и 5, одна 3d-гистограмма для запроса 8.  Каждый запрос был проанализирован с помощью команды EXPLAIN и представлен с помощью графического плана выполнения. Графические планы выполнения были построены в программе dalibo. Также были написаны реляционные формулы выполнения запросов.